\documentclass[10pt,a4paper]{amsart}
\usepackage[margin=19mm]{geometry}
\usepackage{amsmath,amssymb,mathtools}
\usepackage[scr=rsfs,cal=euler]{mathalfa}
\usepackage{xfrac}
\usepackage[unicode,hidelinks,bookmarks=false]{hyperref}
\newcommand{\ie}{i.e.\ }
\newcommand{\cf}{cf.\ }
\newcommand{\resp}{respectively\ }
\newcommand{\Subsection}{\subsection}
\providecommand{\nobreakdash}{\nobreak\hbox{-}}
\setlength{\emergencystretch}{2em}
\multlinegap0pt
\usepackage{amssymb}
\usepackage{enumerate}
\usepackage{xcolor}
\usepackage{tikz}
\usepackage{bm}
\usepackage{float}
\usepackage[normalem]{ulem}
\newtheorem{thm}{Theorem}
\newtheorem{prop}{Proposition}
\newcommand{\N}{\mathbb{N}}
\title{Division properties of commuting polynomials}
\author{Kimiko Hasegawa, Rin Sugiyama}
\date{Source: arXiv:2404.01605v2 (CC BY 4.0)}
\begin{document}
\maketitle

\noindent\emph{Source and licence.} Division properties of commuting polynomials. Version arXiv:2404.01605v2 (CC BY 4.0), distributed by arXiv under the Creative Commons Attribution 4.0 International licence, http://creativecommons.org/licenses/by/4.0/, as recorded in the arXiv licence metadata for this exact version and shown on the abstract page https://arxiv.org/abs/2404.01605v2. Full licence notice: \texttt{licenses/arxiv-2404.01605-CC-BY-4.0.txt}.\par\smallskip

\noindent\textit{Editorial scope.} Counted unit: the article's thm thm-A (source0.tex lines 1287--1292). The article's complete original proof of it is delivered with it (source0.tex lines 1389--1459, one \texttt{proof} environment of 71 lines, which the article places away from the statement). No mathematics is written, repaired or re-derived: the statement and the proof are the article's own lines, in the article's own notation, and no retained line is substituted or re-flowed.  Retained uncounted as the source-local context the proof needs: lines 228--233; lines 1461--1463.  Nothing was omitted from any retained span, so the package carries no omission record.  No retained line contains a citation, so the wrapper carries no bibliography and no bibliography entry.  No retained line contains a figure environment, an image-inclusion command or drawing code, so no image is delivered and the package declares no asset dependency.  Editorial formatting only, in addition: an AMS \texttt{amsart} preamble that restates the article's own declarations and macros used by the retained lines, namely the environments \texttt{thm}, \texttt{prop} on separate counters, together with the standard packages those lines need.  The retained environments print in this document as Theorem 1, Proposition 1 in order of appearance, whereas the article prints the same items with the numbering of its own full document.  The complete native source of the article as distributed in the arXiv e-print for this exact version is bundled unchanged as \texttt{sources/157069/source0.tex}.
\begin{thm}[Theorem \ref{thm-A}]\label{intro-thm-A}
%Assume that $p\geq 3$.
Let $\{f_n(x)\}_{n \in \N}$ be a chain in $k[x]$.
Then $\{f_n(x)\}_{n \in \N}$ is similar to only one of the two $\{ x^n\}_{n\in\N}$ or $\{G_n(x)\}_{n\in\N}$.
Here $G_n(x):=F_n(x) \pmod p$.
\end{thm}

\begin{prop}\label{prop-A2}
There is at most one monic polynomial of degree $k\geq 1$ that commutes with a given monic polynomial of degree 3.
\end{prop}

\begin{thm}[Theorem \ref{intro-thm-A}]\label{thm-A}
%Assume that $p\geq 3$.
Let $\{f_n(x)\}_{n \in \N}$ be a chain in $k[x]$.
Then $\{f_n(x)\}_{n \in \N}$ is similar to only one of the two $\{ x^n\}_{n\in\N}$ or $\{G_n(x)\}_{n\in\N}$.
Here $G_n(x):=F_n(x) \pmod p$.
\end{thm}

\begin{proof}[Proof of Theorem \ref{thm-A}]
We write $f_2(x)=a_2x^2+a_1x+a_0$.
By a similarity via $\lambda(x)=\frac{x}{a_2}$, we may assume that $f_2(x)=x^2+a_1x+a_0$.
Since $\{f_n(x)\}_{n\in \N}$ is a chain, $f_2(x)$ is commute under composition with
\begin{align*}
f_3(x)&=b_3x^3+b_2x^2+b_1x+b_0.
\end{align*}
Comparing the coefficients of $(f_2\circ f_3)(x)$ and $(f_3\circ f_2)(x)$;
\begin{align*}
(f_2\circ f_3)(x)&=b_3^2x^6+b_2^2x^4+a_1b_3x^3+(b_1^2+a_1b_2)x^2+a_1b_1x+b_0^2+a_1b_0+a_0,\\
(f_3\circ f_2)(x)&=b_3x^6+b_3a_1x^5+(b_3a_0+b_3a_1^2+b_2)x^4+b_3a_1^3x^3\\
&\qquad +(b_3a_1^2a_0+b_3a_0^2+b_2a_1^2+b_1)x^2+(b_3a_0^2a_1+b_1a_1)x+b_3a_0^3\\
&\qquad \qquad  +b_2a_0^2+b_1a_0+b_0,
\end{align*}
we see that 
\[
b_3=1,\ a_1=0,\ b_2^2+b_2=a_0,\ b_1^2+b_1=a_0^2,\ b_0^2+a_0=f_3(a_0).
\]
Thus we have $f_2(x)=x^2+b_2^2+b_2=(x+b_2)^2+b_2$.
By a similarity via $\lambda(x)=x+b_2$, we may assume that $f_2(x)=x^2$ and $f_3(x)=x^3+b_1x+b_0$.
In this situation, the equality $(f_2\circ f_3)(x)=(f_3\circ f_2)(x)$ gives
\begin{align}\label{f_3-coef}
b_1^2=b_1,\ b_0^2=b_0.
\end{align}
The similar computation for $f_2(x)=x^2$ and $f_5(x)=c_5x^5+c_4x^4+c_3x^3+c_2x^2+c_1x+c_0$ gives
\begin{align}\label{f_5-coef}
c_5=1,\ c_i^2=c_i\ (i=0,1,2,3,4).
\end{align}

Since $f_3(x)$ also commute under composition with $f_5(x)$, comparing the coefficients of $x^{14}, 
\dots, x^{10}$ in $(f_3\circ f_5)(x)=(f_5\circ f_3)(x)$;
\begin{align}\label{f_3-f_5}
(f_3\circ f_5)(x)
=&x^{15}
+c_4x^{14}
+(c_3+c_4^2)x^{13}
+(c_2+c_4^3)x^{12}
+(c_1+c_3c_4^2+c_3^2)x^{11}\\\notag
&+(c_0+c_2c_4^2+c_3^2c_4)x^{10}
+(c_1c_4^2+c_3^3+c_2^2)x^9%\\\notag
% &+(c_0c_4+b_2c_4+c_2c_3+c_2c_4)x^8
% +(c_1c_3+c_2c_3+c_1)x^7
+\cdots ,\\\notag
(f_5\circ f_3)(x)
=&x^{15}
+b_1x^{13}
+(b_0+c_4)x^{12}
+c_3x^9+\cdots ,
\end{align}
we see that
\begin{align}\label{b-c}
c_4=0,\ c_3=b_1,\ c_2=b_0=0,\ c_1=b_1,\ c_0=0.
\end{align}
Here we used \eqref{f_5-coef}.
By \eqref{f_3-coef}, we now have
\begin{enumerate}
\item 
$b_1=0$ case.
\[
f_3(x)=x^3,\quad f_5(x)=x^5.
\]
\item
$b_1=1$ case.
\begin{align*}
    f_3(x)&=x^3+x=x(x+1)^2=G_3(x),\\ f_5(x)&=x^5+x^3+x=x(x^2+x+1)^2=G_5(x).
\end{align*}

\end{enumerate}
Thus, up to similarity there are two possibilities of chains of degree $\leq 5$.
The assertion now follows from Proposition \ref{prop-A2} below.
\end{proof}

\end{document}
